% s-316-p316-m002; original PDF 316, 317
\Aphorism{s-316-p316-m002}{371}
Die Wissenschaft zerstört nicht die grossen Gefühle (wie die Empfindsamen klagen), sondern nur die verstandesmässige und doch unwissenschaftliche Ausdeutung grosser Gefühle, also jene überflüssige und entstellende Halbheit, auf die gerade der volle gesättigte Gefühlsmensch verzichten kann und gern verzichtet. Sei mitleidig oder liebe Musik — die Wissenschaft wehrt dir’s nicht; aber wenn du sagst „Mitleid durchbricht die Schranken der Erscheinung“ oder „Musik ist unmittelbares Abbild des Willens zum Leben“, so bist du vom reinen Gefühl abgekommen, ohne Erkenntniss dafür zu ernten, und die Wissenschaft schickt dich heim. Gefühl ist das Unwissenschaftliche, Ausserwissenschaftliche, aber nicht Gegenwissenschaftliche; der Kritik der Wissenschaft unterliegt immer nur Wissenschaft — aber auch mit Gefühl verunreinigte Wissenschaft, wissenschaftlich dilettirendes Gefühl. Der Nachempfinder und Mischling wehrt sich gegen die Wissenschaft; der stark und innerlich Fühlende erträgt oder ignorirt sie.
\par

